% !TeX root = main.tex
\section{结论}\label{sec:conclusion}

给定偏好下的多目标序贯决策，关键在于把目标之间的取舍转化为明确的长期策略比较准则，并在求解和评价中始终保持这一准则的含义。本文围绕“长期准则---信息表示---策略类别---模型与数据条件”四个维度，对这一问题的建模、求解与评价研究进行了梳理和比较，主要结论如下。

第一，偏好已知只是问题定义的起点。同一组偏好参数，会因评价对象（阶段结果、累计结果还是轨迹特征）、时间聚合方式、随机结果按 SER 还是 ESR 评价，以及是否附带优先关系和可接受性要求的不同，形成不同的长期准则。SER 与 ESR 的运算顺序决定了评价只需要均值信息，还是需要每条轨迹上的联合结果；联合达标概率和尾部风险与 ESR 一样依赖轨迹层面的信息。用惩罚项代替硬约束、用有限权重近似严格优先关系，都可能改变原准则。

第二，准则的结构决定求解时必须保留的信息。固定线性权重、共同时间结构下的期望可加回报，在满足状态与可行性条件时可以转化为标准的标量决策目标；累计结果上的非线性效用需要通过扩展状态恢复递推结构，代价是更大的状态空间以及更高的计算、数据覆盖和部署信息需求。可递推性与动态一致性是两个不同的要求。最优解是否需要随机化、记忆或时间依赖，同样由准则决定，应作为问题定义的一部分事先明确。

第三，模型可用、允许交互和只有固定历史数据三类信息条件，给出了不同的求解路线和保证范围。模型可用时，可以通过扩展状态规划、逐层动作限制或模型检验查询直接处理非线性效用、优先级和概率要求；允许交互时，非线性 SER 的梯度方法与占用测度方法、非线性 ESR 的分布价值方法以及词典序学习，各自依赖不同的结构条件；只有固定数据时，数据能否重建准则、能否覆盖目标策略的决策区域，是两项互不替代的前提。跨偏好的策略库与条件策略是复用计算与数据的手段，其保证限于训练覆盖的偏好类别和参数区域。

第四，策略评价要回到原准则。单策略评价按原准则的聚合顺序估计性能并报告约束表现；策略集合与条件策略需要把解集质量与输入给定偏好后实际得到的策略表现分开评价；离线评价方法的选择由准则所需的信息决定，均值估计、标量分布估计和多变量联合分布估计各有适用范围。基准环境、仿真实验、历史数据分析与真实系统部署提供的证据层次不同，应用结论的适用范围应与之对应。

在此基础上，本文提出了四个研究方向：复杂长期准则的充分表示与近似误差控制，跨偏好复用的表达范围与性能保证，固定数据下的偏好相关覆盖与可信评价，以及部署信息与模型不确定性下的闭环保证。这四个方向的共同点在于：把表示、估计、优化和执行各环节的误差都落实到原始长期准则上，从而明确有限数据和部署信息能够支持的决策范围。只有在问题定义一致、性能依据可核查的前提下，不同方法的比较和应用验证才有意义。
